\documentclass{ximera}



% MATH -----------------------------------------------------------
\newcommand{\nats}{\mathbb N}
\newcommand{\ints}{\mathbb Z}
\newcommand{\rats}{\mathbb Q}
\newcommand{\reals}{\mathbb R}
\newcommand{\complex}{\mathbb C}
\newcommand{\powerset}{\mathscr P}

%-------------------------------------------------------------

\pgfplotsset{compat=1.18}
\begin{document}
\begin{abstract}
 \end{abstract}

    \title{Class Discussion Activity 9}
    
    \maketitle

\section*{Exercises}

For exercises 1 and 2 use the following information:

A $1$-kg rock is dropped from the top of a large vertical cliff and during free-fall the air resists the motion of the rock proportional to its speed, at a rate of $4$ N of force per m/s of speed.


\begin{enumerate}[label=(\arabic*)]

\item  Determine the IVP that models the vertical velocity $v(t)$ of the rock at $t\geq 0$ seconds after it is dropped.  Use $g=9.8$ m/s$^2$ for the magnitude of the acceleration due to gravity.

\begin{enumerate}[label=(\alph*)]
    \item $v^{\,\prime}=9.8-4v$, $v(0)=0$
    \item $v^{\,\prime}=9.8+4v$, $v(0)=0$
    \item $v^{\,\prime}=-9.8+4v$, $v(0)=0$
    \item $v^{\,\prime}=-9.8-4v$, $v(0)=0$ % ***
    \item $v^{\,\prime}=9.8v+4$, $v(0)=-1$
    \item $v^{\,\prime}=-9.8v-4$, $v(0)=-1$

\end{enumerate}

% (d) ... (f)

\item Given the rock hits the ground at a speed of $2.42$ meters per second determine the height of the cliff.  Round to the nearest meter.

% 2

% -2.42 = -2.45+2.45e^(-4t) ... t_h = -ln(1-2.42/2.45)/4

% Integrate to get 2.45/4*(ln(1-2.42/2.45)+2.42/2.45)\approx -2

\item A $95$ kg skydiver falls through air that resists motion at a rate of $18$ N of force per m/s of speed.  Find the terminal velocity in meters per second.  Use $g=9.8$ m/s$^2$ for the magnitude of the acceleration due to gravity.  Round to the nearest tenth of a meter per second.

% v_infty = -mg/k approx -95*9.8/18 = -51.7 m/s.

\newpage


\item  A $10$ kg object falls attached to vertical rails with a friction-breaking device that is designed to exert a resistive force \emph{proportional to the fourth power of the speed} of the object.  The resistance is $96$ N if the speed is $2$ meter per second.  Find the terminal velocity in meters per second to the nearest hundredth of a meter per second.  Use $g=9.8$ m/s$^2$ for the magnitude of the acceleration due to gravity.

    % -2.01

    % F_d=kv^4 where 96=k2^4 so k=6.

    % 10v'=-98+6v^4=0 so -(98/6)^(1/4)\approx  -2.01 m/s


\item A $1$ kg stone is sling-shot straight up into the air at $40$ meter per second \emph{from an initial height of $1$ meter}.  Assuming the air resists the stone's motion at a rate of $1/25$ N of force per m/s of speed, determine the height in meters reached by the object.  Use $g=9.8$ m/s$^2$ for the magnitude of the acceleration due to gravity.  Round to the nearest tenth of a meter.

% 9.8/(1/25)=245

% v(t)=-245+(40+245)e^{-t/25}.  Set to zero.  b=-25ln(245/285)\approx 3.78

% integral of v(t) from 0 to b is approx. 73.7.  ADD 1 METER!!

% 74.7 is the final answer

\end{enumerate}


\end{document}